\section{A brief systematic review of AI safety research}
\label{sec:slr}

The experiments were motivated by a systematic literature review covering Scopus, Web of Science, IEEE, ACM, and EBSCO, supplemented by snowballing \citep{wohlin_guidelines_2014}. The database search combined terms for artificial intelligence, language models, generative AI, and AGI with ``AI safety.'' It excluded selected application-specific terms. Database eligibility criteria included English-language, peer-reviewed work published from 2024 onward; snowballing additionally incorporated relevant earlier work and gray literature. The final corpus comprised 63 database-search entries and 17 snowballing entries.

The 80 included entries were mapped to seven categories informed by technical safety agendas and AI risk taxonomies \citep{hendrycks2021unsolved,Slattery_2026}. Entries could belong to multiple categories, producing 135 category assignments. Human and societal impacts, value alignment, and governance received the most assignments. Five raters also assessed category severity across magnitude, certainty, irreversibility, and scope. A descriptive comparison with a severity-proportional allocation suggested relatively low coverage of liability and ethics, bias and fairness, and explainability and interpretability. However, ordinal Krippendorff's $\alpha$ ranged from $-0.113$ to $0.122$ across the dimensions. The severity-based comparison therefore does not establish a consensus ranking of research priorities.

Within this corpus, alignment faking emerged as a specific empirical question for smaller models, motivating the following experiments. This choice does not mean that the review established alignment faking as the most severe or most neglected safety problem. The review supplies context for the experimental question; it is not used to infer the prevalence of deception. Category counts and severity summaries are included in Appendix~\ref{app:slr}.

